\scanpage{27}

\[\operatorname{Hom}(\Psi^i_{\unsure{\ast|\bullet}},\Phi(j,n))\longrightarrow
\underset{\substack{\Vert\\0\ \text{si }j\ne\unsure{j|i}}}{\operatorname{Hom}(\Psi^i,\Psi^j)}\]

donc si $j\ne i$, \inserted{$\forall\,\alpha\in$} $\operatorname{Hom}(\Psi^i_{\deleted{\unsure{\ast}}},\Phi(j,n))\xleftarrow{\sim}\operatorname{Hom}(\Psi^i,\Psi^{j-1}\deleted{\unsure{\ne i}})$, \deleted{\unsure{d'où}} qui est nul si $i\ne j-1$ i.e. $j\ne i+1$ ; si $j=i+1$, on trouve $\simeq k$ \unsure{---}.

Pour $j=i$, le composé $\Psi^i\xrightarrow{\alpha}\Phi(i,n)\xrightarrow{\mathrm{can}}\Psi^i$

\diagram{\begin{tikzcd}[ampersand replacement=\&,column sep=small,row sep=small]
\& \& \Psi^i\arrow[d,"\alpha"']\arrow[dr] \& \&\\
0\arrow[r] \& \Psi^{j-1}\arrow[r,dash] \& \Phi(j,n)\arrow[r,dash] \& \Psi^j\arrow[r,dash] \& 0
\end{tikzcd}}

est de la forme $d\cdot\mathrm{id}_{\Psi^i}$, donc l'image inverse \deleted{\unsure{par|p}}de l'extension $\Phi(i,n)$ est de classe $d e_{i,n}=dn e_i$, qui est nulle, donc $dn e_i=0$, donc $dn=0$, donc $d=0$, donc \deleted{\unsure{$\Phi$}} $\alpha$ se factorise par $\Psi^i\mathrel{-}\Psi^{i-1}$, donc est nul.

\deleted{Pour $j=i+1$, on a $\operatorname{Hom}($} Donc \deleted{$\unsure{\operatorname{Hom}(\Phi,\Psi}$}

\[\operatorname{Hom}(\Psi^i,\Phi(j,n))\simeq
\begin{cases}0\ \text{si }j\ne i+1\\ k\ \text{si }j=i+1\ (\text{base \unsure{formé|formée} de}\\
\quad\text{l'hom can }\Psi^i\mathrel{-}\Phi(i+1,n))\end{cases}\]

\noindent\rule{.64\linewidth}{.3pt}

Enfin on suppose

\begin{tabular}{|p{.85\linewidth}}\textcircled{d} Tout objet de $C$ est somme finie d'objets de la forme $\Phi(i,n)$ \inserted{$(\unsure{i|n}\geq1)$} ou $\Psi^j$ $(j\geq0)$\end{tabular}

% archive sheet number: 26
